\section{Related Work}
\label{sec:related}

\paragraph{Activation functions.}
The sigmoid and tanh activations of early multilayer perceptrons saturate, and
their gradients vanish exponentially with depth, which is the classical
explanation for why deep stacks were impractical before the rectifier era.
ReLU removes saturation on the positive side and yields sparse activations;
its variants (LeakyReLU, ELU, SELU, PReLU) address the dying-unit and
negative-saturation issues at some cost in simplicity. Smooth alternatives
that interpolate between the identity on the positive side and a soft
rectifier on the negative side --- GELU and SiLU --- are now the default in
most large-scale vision and language models, and are the activations our
polynomial most directly competes with. We therefore restrict the pointwise
comparison to this deployed set (ReLU, GELU, SiLU) plus tanh and the identity
control, rather than re-auditing the historical variants.

\paragraph{Polynomial and learnable activations.}
Polynomial activations have a long history: the higher-order units of early
connectionist models, and more recently the observation that a polynomial
network of low degree can compute the \emph{exact} product of its inputs,
which a ReLU network can only approximate with growing width. Learnable
activation functions --- PReLU and its family, ``maxout'' networks, and the
recent Kolmogorov--Arnold network (KAN) whose edges carry learnable spline
univariates --- move the nonlinearity into the parameter set. Our study
includes both ends of this axis: a \emph{fixed} cubic, and \emph{learnable}
polynomials in the Legendre, Chebyshev and cubic B-spline bases.

\paragraph{Gated linear units.}
The gated linear unit multiplies a linear projection of the input by a
gate-shaped version of it. Substituting SiLU for the sigmoid in this
construction yields SwiGLU, which is standard in current transformer
feed-forward blocks; GeGLU is the analogous construction with a GELU gate. A
recurring methodological problem in this literature is that gated units carry
roughly $3/2$ times the parameters of the ungated layer they replace, so a
gain in accuracy can reflect capacity rather than gating. We therefore match
parameter count explicitly and report an unmatched-width control.

\paragraph{Normalization layers.}
BatchNorm normalizes each feature using mini-batch statistics, coupling
samples and introducing a train/eval discrepancy resolved by running
statistics. LayerNorm normalizes per sample instead, removing the
batch coupling; RMSNorm drops the mean-centering and is standard in current
transformers. Weight normalization and the $\ell_2$ normalization of weight
vectors are the parameter-free relatives of this family. The $\ell_2$
projection in Eq.~\eqref{eq:projective-poly} belongs to that parameter-free
subfamily, and the question we study --- whether activation \emph{projection}
can substitute for statistical normalization --- is the natural one to ask of
it.

\paragraph{Geometry of deep networks.}
A parallel literature studies deep networks through the geometry of their
loss landscape: the role of saddle points in the training of linear and
residual networks, the projective geometry of the critical locus of a linear
network's gradient equations, and the way normalization and residual
connections reshape that locus. Our construction is the architectural analogue
of that reading --- identifying a sample's activation with a point of
projective space --- and our experimental question, which component of the
construction carries the benefit, is the empirical counterpart of asking which
component of the geometry is responsible for trainability. We note explicitly
that we do \emph{not} claim a theorem here: we treat the projective view as
motivation for the architecture, and we restrict our claims to the measured
behavior of the network.

\paragraph{Empirical comparisons and reproducibility.}
Ablation studies of activations are notoriously sensitive to uncontrolled
factors: parameter count, initialization, learning rate, and depth interact
with the activation choice, and a single-seed comparison at one depth is weak
evidence. Our protocol varies each factor independently, holds the
initialization and optimizer fixed, matches parameter count across
architecturally different families, and repeats the headline comparison over
several random seeds. We also report per-configuration wall-clock cost, since
one of our findings (the poor conditioning of the unnormalized gated unit)
is visible only in the cost profile.
